% Specimen: English thesis in one column, without a bibliography.
\documentclass[preset=thesis, language=english]{.labforge}

\newcommand{\filler}{The arc chamber is the heart of the reactor and the first stage of
the process. It holds the electric arc and provides the space in which nitrogen
dissociates and nitric oxide forms. Its design set most of the geometry of the
apparatus and the requirements of the other parts.}

\begin{document}

\mythesis[title = Fixing atmospheric nitrogen with an electric arc, author = Group B]

\mysection[label=method]{Method}
\filler{} The setup is shown in \myref{setup}.

\myfigure[label=setup, caption=The chamber before assembly., width=0.5]{example-image-a.png}

\filler{} Nitric oxide forms by \myref{formation}; its equilibrium follows
\myref{equilibrium}.

\myreaction[label=formation]{N2 + O2 <=> 2 NO}
\myformula[label=equilibrium]{K_p &= \frac{p_{\mathrm{NO}}^2}{p_{\mathrm{N_2}}\, p_{\mathrm{O_2}}} \\
                              \Delta G^\circ &= -RT \ln K_p}

\myfigure[label=detail, caption=The upper flange., width=0.35, place=wrap]{example-image-b.png}
\filler{} \filler{} The flange is shown in \myref{detail}.

\myfigure[caption=The base mount., width=0.45, place=wrap-left]{example-image-a.png}
\filler{} \filler{}

\mysubsection{Measurements}
\mytable[caption={Mass, reaction time and yield per sample.}]{yields.tex}

\filler{} See \myref{tab:yields, setup} and \myref{method}.
Mixed references: \myref{tab:yields, formation}. Indexed and mixed: \myref{equilibrium[1], formation}.

\end{document}
